% ====================================================================
% РАЗДЕЛ: МЕХАНИКА (СБОРНИК ЗАДАНИЙ №5 И №6)
% ====================================================================

% ====================================================================
% ВАРИАНТ 5
% ====================================================================

% === ЗАДАЧА 5 ===
% Тег: #МЕХАНИКА #КИМ_05 #ВАРИАНТ_05
\begin{taskbasic}[code={\label{task:dem26_v5_n5}}]
\textbf{Задача 5.} Два тела, каждое массой $2$~кг, движутся по одной прямой, вдоль которой направлена ось $Ox$. На рисунке приведены графики зависимости проекций $v_x$ их скоростей от времени $t$. 

\nopagebreak\smallskip
\begin{center}
\begin{tikzpicture}[x=0.25cm, y=0.15cm, every node/.style={font=\footnotesize}]
    \draw[xstep=2, ystep=2, gray!30, thin] (0,0) grid (18,8);
    \draw[-{Stealth[scale=0.8]}, black, thick] (0,0) -- (20,0) node[right] {$t, \text{ с}$};
    \draw[-{Stealth[scale=0.8]}, black, thick] (0,0) -- (0,9) node[above] {$v_x, \text{ м/с}$};
    \node[left] at (0,8) {$8$};
    \node[left] at (0,6) {$6$};
    \node[left] at (0,4) {$4$};
    \node[left] at (0,2) {$2$};
    \node[below left] at (0,0) {$0$};
    \foreach \x in {2,4,6,8,10,12,14,16,18} \node[below] at (\x,0) {\x};
    \draw[ultra thick, brandPrimary] (0,7) -- (18,3.14) node[right, black] {1};
    \draw[ultra thick, brandAccent] (0,0.5) -- (18,5) node[right, black] {2};
\end{tikzpicture}
\end{center}

Из приведённого ниже списка выберите все верные утверждения о движении тел.
1) Равнодействующая сил, действующих на тело 1, монотонно убывает в течение всего времени наблюдения.
2) В промежутке времени от $0$ до $14$~с тела двигались навстречу друг другу.
3) Кинетическая энергия тела 2 за промежуток времени от $6$ до $14$~с увеличилась в $4$ раза.
4) Путь, пройденный телом 1 за промежуток времени от $0$ до $18$~с, больше пути, пройденного телом 2 за тот же промежуток времени.
5) Модуль импульса тела 2 в момент времени $t=10$~с равен $10$~\text{кг}$\cdot$\text{м/с}.

\kim{5}
\end{taskbasic}
\answer{34}


% ====================================================================
% ВАРИАНТ 7
% ====================================================================

% === ЗАДАЧА 5 ===
% Тег: #МЕХАНИКА #КИМ_05 #ВАРИАНТ_07
\begin{taskbasic}[code={\label{task:dem26_v7_n5}}]
\textbf{Задача 5.} На рисунке показан график зависимости координаты $x$ тела, движущегося вдоль оси $Ox$, от времени $t$. 

\nopagebreak\smallskip
\begin{center}
\begin{tikzpicture}[scale=1.2, every node/.style={font=\footnotesize}]
    \draw[-{Stealth[scale=0.8]}, black, thick] (0,0) -- (0,2.5) node[above] {$x$};
    \draw[-{Stealth[scale=0.8]}, black, thick] (0,0) -- (4,0) node[right] {$t$};
    \draw[ultra thick, brandPrimary, domain=0.5:3.5, samples=100] plot (\x, {1.2 + sin(120*(\x-1.5))});
    \fill[black] (0.8, 0.45) circle (1.5pt) node[left] {$A$};
    \fill[black] (1.15, 0.2) circle (1.5pt) node[below] {$B$};
    \fill[black] (2.25, 1.95) circle (1.5pt) node[left] {$C$};
    \fill[black] (3.1, 1.45) circle (1.5pt) node[right] {$D$};
    \node[below left] at (0,0) {0};
\end{tikzpicture}
\end{center}

Из приведённого ниже списка выберите все верные утверждения.
1) В положении $A$ модуль скорости тела больше, чем в положении $D$.
2) В точке $B$ проекция скорости тела на ось $Ox$ равна нулю.
3) В положении $D$ векторы скорости и ускорения тела направлены в противоположные стороны.
4) Проекция перемещения тела на ось $Ox$ при его переходе из точки $A$ в точку $B$ положительна.
5) На участке $AB$ модуль скорости тела уменьшается.

\kim{5}
\end{taskbasic}
\answer{125}


% ====================================================================
% ВАРИАНТ 13
% ====================================================================

% === ЗАДАЧА 5 ===
% Тег: #МЕХАНИКА #КИМ_05 #ВАРИАНТ_13
\begin{taskbasic}[code={\label{task:dem26_v13_n5}}]
\textbf{Задача 5.} На рисунке показан график зависимости координаты $x$ тела, движущегося вдоль оси $Ox$, от времени $t$. 

\nopagebreak\smallskip
\begin{center}
\begin{tikzpicture}[scale=1.2, every node/.style={font=\footnotesize}]
    \draw[-{Stealth[scale=0.8]}, black, thick] (0,0) -- (0,2.5) node[above] {$x$};
    \draw[-{Stealth[scale=0.8]}, black, thick] (0,0) -- (4,0) node[right] {$t$};
    \draw[ultra thick, brandPrimary, domain=0.5:3.5, samples=100] plot (\x, {1.2 + sin(120*(\x-1.5))});
    \fill[black] (0.8, 0.45) circle (1.5pt) node[left] {$A$};
    \fill[black] (1.9, 1.55) circle (1.5pt) node[left] {$B$};
    \fill[black] (2.65, 2.2) circle (1.5pt) node[above] {$C$};
    \fill[black] (3.3, 1.1) circle (1.5pt) node[right] {$D$};
    \node[below left] at (0,0) {0};
\end{tikzpicture}
\end{center}

Из приведённого ниже списка выберите все верные утверждения.
1) В точке $A$ ускорение тела равно нулю.
2) В точке $B$ проекция ускорения тела на ось $Ox$ положительна.
3) Проекция перемещения тела на ось $Ox$ при переходе из точки $B$ в точку $C$ отрицательна.
4) В точке $D$ проекция скорости тела на ось $Ox$ отрицательна.
5) На участке $CD$ кинетическая энергия тела увеличивается.

\kim{5}
\end{taskbasic}
\answer{13}


% ====================================================================
% ВАРИАНТ 15
% ====================================================================

% === ЗАДАЧА 6 ===
% Тег: #МЕХАНИКА #КИМ_06 #ВАРИАНТ_15
\begin{taskbasic}[code={\label{task:dem26_v15_n6}}]
\textbf{Задача 6.} На рисунке показан график зависимости координаты $x$ тела, движущегося вдоль оси $Ox$, от времени $t$ (парабола). Графики А и Б представляют собой зависимости физических величин, характеризующих движение этого тела, от времени $t$.

\nopagebreak\smallskip
\begin{center}
\begin{minipage}[b]{0.3\linewidth}
\centering
\begin{tikzpicture}[scale=0.6, every node/.style={font=\footnotesize}]
    \draw[-{Stealth[scale=0.8]}, black, thick] (0,0) -- (0,2.5) node[above] {$x$};
    \draw[-{Stealth[scale=0.8]}, black, thick] (0,0) -- (3,0) node[right] {$t$};
    \draw[ultra thick, brandPrimary, domain=0:2.5] plot (\x, {1.5 - 1.2*\x + 0.3*\x*\x});
    \node[below] at (2,0) {$t_1$};
    \node[below left] at (0,0) {0};
\end{tikzpicture}
\end{minipage}
\begin{minipage}[b]{0.3\linewidth}
\centering
\begin{tikzpicture}[scale=0.6, every node/.style={font=\footnotesize}]
    \node[above] at (1.2,2.0) {А)};
    \draw[-{Stealth[scale=0.8]}, black, thick] (0,0) -- (0,2.2) node[above] {};
    \draw[-{Stealth[scale=0.8]}, black, thick] (0,0) -- (3,0) node[right] {$t$};
    \draw[ultra thick, brandPrimary] (0,-1) -- (2,0) -- (3,0.5);
    \node[below] at (2,0) {$t_1$};
    \node[below left] at (0,0) {0};
\end{tikzpicture}
\end{minipage}
\begin{minipage}[b]{0.3\linewidth}
\centering
\begin{tikzpicture}[scale=0.6, every node/.style={font=\footnotesize}]
    \node[above] at (1.2,2.0) {Б)};
    \draw[-{Stealth[scale=0.8]}, black, thick] (0,0) -- (0,2.2) node[above] {};
    \draw[-{Stealth[scale=0.8]}, black, thick] (0,0) -- (3,0) node[right] {$t$};
    \draw[ultra thick, brandPrimary] (0,0.8) -- (3,0.8);
    \node[below left] at (0,0) {0};
\end{tikzpicture}
\end{minipage}
\end{center}

Установите соответствие между графиками и физическими величинами, зависимость которых от времени эти графики могут представлять.
1) проекция скорости тела на ось $Ox$
2) проекция перемещения тела на ось $Ox$
3) кинетическая энергия тела
4) модуль  равнодействующей сил, действующих на тело

Запишите в таблицу выбранные цифры под соответствующими буквами.

\kim{6}
\end{taskbasic}
\answer{14}


% ====================================================================
% ВАРИАНТ 17
% ====================================================================

% === ЗАДАЧА 6 ===
% Тег: #МЕХАНИКА #КИМ_06 #ВАРИАНТ_17
\begin{taskbasic}[code={\label{task:dem26_v17_n6}}]
\textbf{Задача 6.} После удара в момент времени $t=0$ шайба начала скользить вверх по гладкой наклонной плоскости с начальной скоростью $\vec{v}_0$, как показано на рисунке. В момент времени $t_0$ шайба вернулась в исходное положение. Графики А и Б отображают изменение с течением времени физических величин, характеризующих движение шайбы.

\nopagebreak\smallskip
\begin{center}
\begin{minipage}[b]{0.3\linewidth}
\centering
\begin{tikzpicture}[scale=0.7, every node/.style={font=\footnotesize}]
    \draw[thick] (0,0) -- (3,1.2) -- (3,0) -- cycle;
    \draw[fill=gray!60] (0.8,0.32) rectangle (1.3,0.7) [rotate around={21.8:(1.05,0.51)}];
    \draw[-{Stealth[scale=0.8]}, ultra thick, black] (1.05,0.51) -- (1.8,0.81) node[above] {$\vec{v}_0$};
    \draw[->] (0,0) -- (3.5,0) node[right] {$x$};
    \draw[->] (0,0) -- (0,2.0) node[above] {$y$};
    \node[below left] at (0,0) {0};
\end{tikzpicture}
\end{minipage}
\begin{minipage}[b]{0.3\linewidth}
\centering
\begin{tikzpicture}[scale=0.6, every node/.style={font=\footnotesize}]
    \node[above] at (1.2,2.0) {А)};
    \draw[-{Stealth[scale=0.8]}, black, thick] (0,0) -- (0,2.2) node[above] {};
    \draw[-{Stealth[scale=0.8]}, black, thick] (0,0) -- (3,0) node[right] {$t$};
    \draw[ultra thick, brandPrimary] (0,1.5) -- (3,1.5);
    \node[below] at (3,0) {$t_0$};
    \node[below left] at (0,0) {0};
\end{tikzpicture}
\end{minipage}
\begin{minipage}[b]{0.3\linewidth}
\centering
\begin{tikzpicture}[scale=0.6, every node/.style={font=\footnotesize}]
    \node[above] at (1.2,2.0) {Б)};
    \draw[-{Stealth[scale=0.8]}, black, thick] (0,-1.2) -- (0,1.5) node[above] {};
    \draw[-{Stealth[scale=0.8]}, black, thick] (0,0) -- (3,0) node[right] {$t$};
    \draw[ultra thick, brandPrimary] (0,1.2) -- (3,-1.0);
    \node[below] at (3,0) {$t_0$};
    \node[below left] at (0,0) {0};
\end{tikzpicture}
\end{minipage}
\end{center}

Установите соответствие между графиками и физическими величинами, изменение которых со временем эти графики могут отображать.
1) проекция скорости $v_y$
2) проекция ускорения $a_x$
3) кинетическая энергия $E_{\text{к}}$
4) полная механическая энергия $E_{\text{мех}}$

Запишите в таблицу выбранные цифры под соответствующими буквами.

\kim{6}
\end{taskbasic}
\answer{41}

% === ЗАДАЧА 22 ===
% Тег: #МЕХАНИКА #КИМ_22 #ВАРИАНТ_17
\begin{taskbasic}[code={\label{task:dem26_v17_n22}}]
\textbf{Задача 22.} Прямолинейно движущийся поезд последние перед остановкой $1200$~м своего пути преодолел за $2$~минуты. Определите скорость поезда за $30$~с до остановки. Ускорение поезда считать постоянным.

\kim{22}
\end{taskbasic}
\answer{10~м/с}


% ====================================================================
% ВАРИАНТ 19
% ====================================================================

% === ЗАДАЧА 6 ===
% Тег: #МЕХАНИКА #КИМ_06 #ВАРИАНТ_19
\begin{taskbasic}[code={\label{task:dem26_v19_n6}}]
\textbf{Задача 6.} Тело движется вдоль оси $Ox$, при этом его координата изменяется с течением времени в соответствии с уравнением $x(t)=8+3t-5t^{2}$ (все величины выражены в СИ). Графики А и Б представляют собой зависимости физических величин, характеризующих движение этого тела, от времени $t$.

\nopagebreak\smallskip
\begin{center}
\begin{minipage}[b]{0.45\linewidth}
\centering
\begin{tikzpicture}[scale=0.7, every node/.style={font=\footnotesize}]
    \node[above] at (1.2,2.0) {А)};
    \draw[-{Stealth[scale=0.8]}, black, thick] (0,0) -- (0,2.2) node[above] {};
    \draw[-{Stealth[scale=0.8]}, black, thick] (0,0) -- (3,0) node[right] {$t$};
    \draw[ultra thick, brandPrimary, domain=0:2.5, samples=100] plot (\x, {1.5 - 1.5*\x + 0.5*\x*\x});
    \node[below] at (1.5,0) {$t_1$};
    \node[below left] at (0,0) {0};
\end{tikzpicture}
\end{minipage}
\begin{minipage}[b]{0.45\linewidth}
\centering
\begin{tikzpicture}[scale=0.7, every node/.style={font=\footnotesize}]
    \node[above] at (1.2,2.0) {Б)};
    \draw[-{Stealth[scale=0.8]}, black, thick] (0,-1.2) -- (0,1.5) node[above] {};
    \draw[-{Stealth[scale=0.8]}, black, thick] (0,0) -- (3,0) node[right] {$t$};
    \draw[ultra thick, brandPrimary] (0,0.8) -- (2.5,-1.0);
    \node[below] at (1.0,0) {$t_1$};
    \node[below left] at (0,0) {0};
\end{tikzpicture}
\end{minipage}
\end{center}

Установите соответствие между графиками и физическими величинами, зависимость которых от времени эти графики могут представлять.
1) проекция $a_x$ ускорения тела
2) кинетическая энергия тела
3) модуль импульса тела
4) проекция $v_x$ скорости тела

Запишите в таблицу выбранные цифры под соответствующими буквами.

\kim{6}
\end{taskbasic}
\answer{24}


% ====================================================================
% ВАРИАНТ 20
% ====================================================================

% === ЗАДАЧА 6 ===
% Тег: #МЕХАНИКА #КИМ_06 #ВАРИАНТ_20
\begin{taskbasic}[code={\label{task:dem26_v20_n6}}]
\textbf{Задача 6.} Тело движется вдоль оси $Ox$, при этом его координата изменяется с течением времени в соответствии с уравнением $x(t)=10+2t-6t^{2}$ (все величины выражены в СИ). Графики А и Б представляют собой зависимости физических величин, характеризующих движение этого тела, от времени $t$.

\nopagebreak\smallskip
\begin{center}
\begin{minipage}[b]{0.45\linewidth}
\centering
\begin{tikzpicture}[scale=0.7, every node/.style={font=\footnotesize}]
    \node[above] at (1.2,2.0) {А)};
    \draw[-{Stealth[scale=0.8]}, black, thick] (0,-1.5) -- (0,1.5) node[above] {};
    \draw[-{Stealth[scale=0.8]}, black, thick] (0,0) -- (3,0) node[right] {$t$};
    \draw[ultra thick, brandPrimary] (0,-1.0) -- (3,-1.0);
    \node[below] at (1.0,0) {$t_1$};
    \node[below left] at (0,0) {0};
\end{tikzpicture}
\end{minipage}
\begin{minipage}[b]{0.45\linewidth}
\centering
\begin{tikzpicture}[scale=0.7, every node/.style={font=\footnotesize}]
    \node[above] at (1.2,2.0) {Б)};
    \draw[-{Stealth[scale=0.8]}, black, thick] (0,-1.2) -- (0,1.5) node[above] {};
    \draw[-{Stealth[scale=0.8]}, black, thick] (0,0) -- (3,0) node[right] {$t$};
    \draw[ultra thick, brandPrimary] (0,1.0) -- (2.5,-1.2);
    \node[below] at (1.0,0) {$t_1$};
    \node[below left] at (0,0) {0};
\end{tikzpicture}
\end{minipage}
\end{center}

Установите соответствие между графиками и физическими величины, зависимость которых от времени эти графики могут представлять.
1) проекция $a_x$ ускорения тела
2) кинетическая энергия тела
3) модуль импульса тела
4) проекция $v_x$ скорости тела

Запишите в таблицу выбранные цифры под соответствующими буквами.

\kim{6}
\end{taskbasic}
\answer{14}


% ====================================================================
% ВАРИАНТ 27
% ====================================================================

% === ЗАДАЧА 5 ===
% Тег: #МЕХАНИКА #КИМ_05 #ВАРИАНТ_27
\begin{taskbasic}[code={\label{task:dem26_v27_n5}}]
\textbf{Задача 5.} На рисунке приведены графики зависимости проекции $v_x$ скорости от времени $t$ для двух тел А и В, движущихся вдоль оси $Ox$. 

\nopagebreak\smallskip
\begin{center}
\begin{tikzpicture}[x=0.45cm, y=0.1cm, every node/.style={font=\footnotesize}]
    \draw[xstep=2, ystep=10, gray!30, thin] (0,-20) grid (10,30);
    \draw[-{Stealth[scale=0.8]}, black, thick] (0,0) -- (11,0) node[right] {$t, \text{ с}$};
    \draw[-{Stealth[scale=0.8]}, black, thick] (0,-20) -- (0,34) node[above] {$v_x, \text{ м/с}$};
    \node[left] at (0,30) {$30$};
    \node[left] at (0,20) {$20$};
    \node[left] at (0,10) {$10$};
    \node[left] at (0,-10) {$-10$};
    \node[left] at (0,-20) {$-20$};
    \node[below left] at (0,0) {$0$};
    \foreach \x in {2,4,6,8,10} \node[below] at (\x,0) {\x};
    \draw[ultra thick, brandPrimary] (0,-20) -- (10,30) node[right, black] {$A$};
    \draw[ultra thick, brandAccent, domain=0:8, samples=100] plot (\x, {30 - 1.25*(\x-2)*(\x-2)}) node[right, black] {$B$};
\end{tikzpicture}
\end{center}

Выберите все верные утверждения о характере движения тел.
1) Тело А движется равномерно со скоростью $5$~м/с.
2) Тело В за первые $4$~секунды движения переместилось в положительном направлении оси $Ox$ на расстояние более $120$~м.
3) В промежутке времени от $0$ до $8$~с путь, пройденный телом А, равен $0$.
4) В интервале времени от $6$ до $8$~с тело В изменило направление своего движения.
5) В момент времени $6$~с тела А и В двигались с одинаковой скоростью в одном направлении.

\kim{5}
\end{taskbasic}
\answer{5}

% === ЗАДАЧА 6 ===
% Тег: #МЕХАНИКА #КИМ_06 #ВАРИАНТ_27
\begin{taskbasic}[code={\label{task:dem26_v27_n6}}]
\textbf{Задача 6.} После удара в момент $t=0$ шайба начала скользить вверх по гладкой наклонной плоскости с начальной скоростью, как показано на рисунке, и в момент времени $t=t_0$ вернулась в исходное положение. Графики А и Б отображают изменение с течением времени физических величин, характеризующих движение шайбы.

\nopagebreak\smallskip
\begin{center}
\begin{minipage}[b]{0.3\linewidth}
\centering
\begin{tikzpicture}[scale=0.7, every node/.style={font=\footnotesize}]
    \draw[thick] (0,0) -- (3,1.2) -- (3,0) -- cycle;
    \draw[fill=gray!60] (0.8,0.32) rectangle (1.3,0.7) [rotate around={21.8:(1.05,0.51)}];
    \draw[-{Stealth[scale=0.8]}, ultra thick, black] (1.05,0.51) -- (1.8,0.81) node[above] {$\vec{v}_0$};
    \draw[->] (0,0) -- (3.5,0) node[right] {$x$};
    \draw[->] (0,0) -- (0,2.0) node[above] {$y$};
    \node[below left] at (0,0) {0};
\end{tikzpicture}
\end{minipage}
\begin{minipage}[b]{0.3\linewidth}
\centering
\begin{tikzpicture}[scale=0.6, every node/.style={font=\footnotesize}]
    \node[above] at (1.2,2.0) {А)};
    \draw[-{Stealth[scale=0.8]}, black, thick] (0,0) -- (0,2.2) node[above] {};
    \draw[-{Stealth[scale=0.8]}, black, thick] (0,0) -- (3,0) node[right] {$t$};
    \draw[ultra thick, brandPrimary, domain=0:3, samples=100] plot (\x, {1.5 - 1.5*(\x-1.5)*(\x-1.5)});
    \node[below] at (3,0) {$t_0$};
    \node[below left] at (0,0) {0};
\end{tikzpicture}
\end{minipage}
\begin{minipage}[b]{0.3\linewidth}
\centering
\begin{tikzpicture}[scale=0.6, every node/.style={font=\footnotesize}]
    \node[above] at (1.2,2.0) {Б)};
    \draw[-{Stealth[scale=0.8]}, black, thick] (0,-1.2) -- (0,1.5) node[above] {};
    \draw[-{Stealth[scale=0.8]}, black, thick] (0,0) -- (3,0) node[right] {$t$};
    \draw[ultra thick, brandPrimary] (0,1.2) -- (3,-1.2);
    \node[below] at (3,0) {$t_0$};
    \node[below left] at (0,0) {0};
\end{tikzpicture}
\end{minipage}
\end{center}

Установите соответствие между графиками и физическими величинами, изменение которых со временем эти графики могут отображать.
1) кинетическая энергия $E_{\text{к}}$
2) проекция скорости $v_y$
3) координата $x$
4) проекция силы тяжести на ось $Ox$

Запишите в таблицу выбранные цифры под соответствующими буквами.

\kim{6}
\end{taskbasic}
\answer{32}


% ====================================================================
% ВАРИАНТ 29
% ====================================================================

% === ЗАДАЧА 6 ===
% Тег: #МЕХАНИКА #КИМ_06 #ВАРИАНТ_29
\begin{taskbasic}[code={\label{task:dem26_v29_n6}}]
\textbf{Задача 6.} На рисунке показан график зависимости координаты $x$ тела, движущегося вдоль оси $Ox$, от времени $t$ (парабола). Графики А и Б представляют собой зависимости физических величин, характеризующих движение тела, от времени $t$.

\nopagebreak\smallskip
\begin{center}
\begin{minipage}[b]{0.3\linewidth}
\centering
\begin{tikzpicture}[scale=0.6, every node/.style={font=\footnotesize}]
    \draw[-{Stealth[scale=0.8]}, black, thick] (0,-1.5) -- (0,1.5) node[above] {$x$};
    \draw[-{Stealth[scale=0.8]}, black, thick] (0,0) -- (3,0) node[right] {$t$};
    \draw[ultra thick, brandPrimary, domain=0:2.5] plot (\x, {-0.5 - 0.8*\x + 0.4*\x*\x});
    \node[below left] at (0,0) {0};
\end{tikzpicture}
\end{minipage}
\begin{minipage}[b]{0.3\linewidth}
\centering
\begin{tikzpicture}[scale=0.6, every node/.style={font=\footnotesize}]
    \node[above] at (1.2,2.0) {А)};
    \draw[-{Stealth[scale=0.8]}, black, thick] (0,-1.2) -- (0,1.5) node[above] {};
    \draw[-{Stealth[scale=0.8]}, black, thick] (0,0) -- (3,0) node[right] {$t$};
    \draw[ultra thick, brandPrimary] (0,-1.0) -- (1.2,0) -- (3,1.5);
    \node[below left] at (0,0) {0};
\end{tikzpicture}
\end{minipage}
\begin{minipage}[b]{0.3\linewidth}
\centering
\begin{tikzpicture}[scale=0.6, every node/.style={font=\footnotesize}]
    \node[above] at (1.2,2.0) {Б)};
    \draw[-{Stealth[scale=0.8]}, black, thick] (0,0) -- (0,2.2) node[above] {};
    \draw[-{Stealth[scale=0.8]}, black, thick] (0,0) -- (3,0) node[right] {$t$};
    \draw[ultra thick, brandPrimary] (0,1.2) -- (3,1.2);
    \node[below left] at (0,0) {0};
\end{tikzpicture}
\end{minipage}
\end{center}

Установите соответствие между графиками и физическими величинами, зависимость которых от времени эти графики могут представлять.
1) модуль скорости тела
2) модуль ускорения тела
3) кинетическая энергия тела
4) проекция на ось $Ox$ перемещения тела из начального положения

Запишите в таблицу выбранные цифры под соответствующими буквами.

\kim{6}
\end{taskbasic}
\answer{12}